\title{QLoRA: Efficient Finetuning of Quantized LLMs}

\begin{abstract}
We introduce \textsc{QLoRA}, a memory-efficient finetuning technique that enables the finetuning of a 65B parameter language model on a single 48GB GPU, all while maintaining the performance of conventional 16-bit finetuning on downstream tasks. \textsc{QLoRA} achieves this by performing gradient backpropagation through a fixed, 4-bit quantized pretrained language model, directing updates exclusively into Low Rank Adapters~(LoRA). The resulting model family, called \textbf{Guanaco}, surpasses all previously open-sourced models on the Vicuna benchmark, attaining 99.3\% of ChatGPT’s performance after just 24 hours of single-GPU finetuning. \textsc{QLoRA} achieves significant memory savings through several key contributions: (a) 4-bit NormalFloat (NF4), a novel datatype that is theoretically optimal for representing normally distributed weights; (b) Double Quantization, which compresses memory further by quantizing quantization parameters; and (c) Paged Optimizers, designed to mitigate temporary surges in memory usage. Utilizing \textsc{QLoRA}, we finetuned over 1,000 models, systematically evaluating instruction-following and chatbot capabilities across 8 instruction datasets, a variety of model types (LLaMA, T5), and parameter scales (including challenging 33B and 65B models), which are ordinarily impractical with standard finetuning methods. Our findings demonstrate that finetuning with QLoRA on compact, high-quality datasets can achieve state-of-the-art outcomes, even for models with fewer parameters compared to previous leading approaches. Additionally, we thoroughly analyze chatbot performance based on human assessments and GPT-4-based evaluations, showing that the latter can serve as a cost-effective and robust proxy for human evaluation. We also highlight that current benchmarks are unreliable indicators of true chatbot quality. Finally, a “lemon-picked” evaluation illustrates specific shortcomings of \textbf{Guanaco} relative to ChatGPT. All models and code—including CUDA kernels for 4-bit training—are publicly released.\footnote{\url{https://github.com/artidoro/qlora} and \url{https://github.com/TimDettmers/bitsandbytes}}
\end{abstract}

\section{Introduction}

Adapting large language models (LLMs) through finetuning has emerged as a powerful approach to enhance their performance \citep{min2021metaicl, wei2021finetuned, ouyang2022training, wang2022super, wang2022self, liu2022few} and to tailor or correct specific behaviors \citep{ouyang2022training,askell2021general,bai2022training}. Nonetheless, finetuning models at this scale is prohibitively resource-intensive; for example, standard 16-bit finetuning of a LLaMA model with 65 billion parameters~\citep{touvron2023llama} demands in excess of 780 GB of GPU memory. Although recent advancements in quantization have succeeded in reducing the memory consumption of LLMs during inference \citep{dettmers2022llm, dettmers2022case, frantar2022gptq, xiao2022smoothquant}, such techniques typically fail during training \citep{wortsman2023stable}.

In this work, we establish for the first time that one can successfully finetune a model quantized to 4 bits without sacrificing accuracy. Our approach, \textsc{QLoRA}, applies a novel high-precision quantization scheme to pretrained models, converting them to 4-bit representations, and then incorporates a compact set of trainable Low-rank Adapter weights \citep{hu2021lora}, 

which are optimized by propagating gradients through the quantized model weights. 

By employing \textsc{QLoRA}, the memory demands for finetuning a 65B model are reduced dramatically from over 780GB to under 48GB of GPU memory, all without any detriment to efficiency or predictive performance compared to the conventional 16-bit finetuning approach. This innovation greatly democratizes LLM finetuning, rendering it possible for even the largest publicly available models to be adapted using only a single GPU. Leveraging \textsc{QLoRA}, we introduce the \textbf{Guanaco} suite of models; notably, the second largest model achieves 97.8\% of ChatGPT's performance on the Vicuna benchmark~\citep{vicuna2023} and can be finetuned within 12 hours using just one consumer-grade GPU. With a professional GPU over a 24-hour period, our largest model attains 99.3\%, effectively matching ChatGPT's performance on the same benchmark. Furthermore, when deployed, our smallest \textbf{Guanaco} model, containing 7B parameters, operates with only 5 GB of memory and exceeds the performance of a 26 GB Alpaca model by more than 20 percentage points on the Vicuna benchmark (Table~\ref{tbl:vicuna_eval}).

\textsc{QLoRA} incorporates several novel strategies aimed at minimizing memory consumption while maintaining robust performance: (1) \textbf{4-bit NormalFloat}, an information-theoretically optimal quantization format specifically tailored for normally distributed data, demonstrating superior empirical outcomes in comparison to both 4-bit Integer and 4-bit Float representations. (2) \textbf{Double Quantization}, a technique where even the quantization coefficients themselves are further quantized, resulting in average savings of roughly 0.37 bits per parameter (which equates to approximately 3 GB for models with 65B parameters). (3) \textbf{Paged Optimizers}, which leverage NVIDIA unified memory to eliminate the pronounced memory usage spikes typically experienced during gradient checkpointing when handling mini-batches with extended sequence lengths. These advancements are integrated into an enhanced LoRA methodology that positions adapters within every layer of the network, thereby virtually eliminating the accuracy compromises observed in previous methods.

The increased efficiency of \textsc{QLoRA} makes it possible to systematically examine instruction finetuning and chatbot capabilities at model scales unattainable with standard finetuning approaches due to their substantial memory requirements. Consequently, we train over 1,000 models, encompassing diverse instruction tuning datasets, architectural variants, and parameter scales ranging from 80M to 65B. Our results demonstrate that \textsc{QLoRA} attains parity with 16-bit training (\S\ref{sec:comp_to_std}), supports the training of leading-edge chatbots such as \textbf{Guanaco} (\S\ref{sec:pushing}), and allows for a comprehensive investigation of emergent patterns in trained models. Notably, our findings indicate that the quality of training data plays a much greater role than dataset size; for instance, a compact 9k-sample dataset (OASST1) surpasses a much larger 450k-sample dataset (FLAN v2, subsampled) in chatbot task performance, despite both aiming to enhance instruction-following generalization. Furthermore, we establish that high scores on the Massive Multitask Language Understanding (MMLU) benchmark do not necessarily translate into superior performance on the Vicuna chatbot benchmark and vice versa, underscoring that appropriateness of the dataset to the task is more critical than sheer volume.

In addition, we conduct a comprehensive evaluation of chatbot models using both assessments from human annotators and automated judgments from GPT-4. Our evaluation framework employs a tournament-style competition in which chatbots are paired in head-to-head matches to determine the optimal response for a particular prompt. Outcomes of these matches are adjudicated by either GPT-4 or human evaluators. We then aggregate the results of these tournaments into Elo scores~\citep{elo1967proposed, elo1978rating}, providing a performance-based ranking of the chatbots. Our findings indicate a high degree of concordance between human and GPT-4 evaluations regarding model rankings in the tournaments, though there are notable cases of substantial divergence. Therefore, we underscore that while model-based evaluation presents a cost-effective alternative to human annotation, it introduces its own sources of variability and uncertainty.

To supplement the quantitative benchmark results, we conduct an in-depth qualitative investigation of the \textbf{Guanaco} models. This analysis exposes patterns of strengths and limitations that quantitative methods did not fully capture.

All model outputs, including corresponding annotations from both human raters and GPT-4, are publicly released to support further research. Additionally, we provide open access to our full codebase, CUDA kernel implementations, and integrate our methodology with the Hugging Face transformers library~\citep{wolf2019huggingface}, enhancing accessibility for the community. We make available a suite of adapter modules for models of 7/13/33/65B parameter sizes, each trained on eight separate instruction-following datasets, totaling 32 distinct fine-tuned, open-source model variants.

\section{Background}
\paragraph{Block-wise k-bit Quantization} The quantization process involves mapping data from a representation with high informational content to one with lower informational content. This frequently entails converting numerical data types from higher-precision formats, such as 32-bit floating point, to lower-precision formats, such as 8-bit integer. To ensure optimal use of the target bit-width's value range, it is standard practice to normalize the input data (often structured as a tensor) by its absolute maximum before quantization. For instance, to quantize a tensor of 32-bit floats (FP32) to an 8-bit integer (Int8) representation within the interval $[-127, 127]$, we apply:

\begin{equation}
    \mathbf{X}^{ \text{Int8}} = \text{round}\left( \frac{127}{{\text{absmax}}(\mathbf{X}^{{\text{FP32}}})}\mathbf{X}^{{\text{FP32}}} \right) = \text{round}(c^{\text{FP32}}\cdot \mathbf{X}^{{\text{FP32}}}),
\end{equation}

where the coefficient $c$ denotes the {\em quantization constant} or {\em quantization scale}. The corresponding dequantization procedure is the reverse operation:

\begin{equation}
    \text{dequant}(c^{\text{FP32}}, \mathbf{X}^{\text{Int8}}) = \frac{\mathbf{X}^{\text{Int8}}}{c^{\text{FP32}}} = \mathbf{X}^{\text{FP32}}
\end{equation}

A drawback of this method is that, when the input tensor contains a value with extremely large magnitude (i.e., an outlier), the quantization bins—certain bit patterns—are poorly utilized, resulting in few or no values being assigned to some bins. To address this problem, a widely adopted solution is to partition the input tensor into multiple blocks that are each quantized independently, with a separate quantization constant $c$ for each block. This strategy can be described as follows: The input tensor $\mathbf{X}\in \mathbb{R}^{b\times h}$ is divided into $n$ contiguous blocks of length $B$ by first flattening the tensor, then slicing the linearized array into ${n = ({b\times h} )/{B} }$ segments. Each block is then quantized independently, using Equation~1, resulting in a quantized tensor and a set of $n$ quantization constants $c_i$.

\paragraph{Low-rank Adapters} The Low-rank Adapter (LoRA) finetuning approach \citep{hu2021lora} is designed to decrease memory consumption by introducing a compact collection of trainable weights—referred to as adapters—while the core model weights remain unchanged. During stochastic gradient descent, gradients are propagated through the fixed pretrained parameters to the adapter, which is trained to minimize the loss. LoRA modifies the linear projection layer by injecting an extra, factorized projection. For a given projection ${\mathbf{X} \mathbf{W} = \mathbf{Y}}$, where $\mathbf{X}\in  \mathbb{R}^{b\times h}$ and $\mathbf{W}\in  \mathbb{R}^{h\times o}$, the LoRA update is given by:

\begin{equation}
    \mathbf{Y} = \mathbf{X}\mathbf{W} + s\mathbf{X}\mathbf{L}_1\mathbf{L}_2,
\end{equation}

where $\mathbf{L}_1\in  \mathbb{R}^{h\times r}$, $\mathbf{L}_2\in  \mathbb{R}^{r\times o}$, and $s$ denotes a scalar factor.

\paragraph{Memory Requirement of Parameter-Efficient Finetuning} A significant aspect to consider involves the memory demands of LoRA during training, particularly with respect to both the quantity and scale of adapters utilized. Due to the exceptionally small memory consumption associated with LoRA, incorporating a greater number of adapters to enhance performance is possible without a substantial rise in overall memory utilization. Although LoRA is classified as a Parameter Efficient Finetuning (PEFT) technique, in the context of LLM finetuning, the dominant contributor to memory usage stems from activation gradients, rather than the trainable parameters introduced by LoRA. For instance, using a 7B LLaMA model on FLAN v2 with a batch size of 1, and configuring LoRA with weights equaling the frequently adopted 0.2\% of the base model's parameters~\citep{hu2021lora,liu2022few}, the memory consumption of LoRA input gradients is 567 MB, whereas the LoRA parameters themselves only occupy 26 MB. When applying gradient checkpointing~\citep{chen2016training}, the memory usage for input gradients drops to an average of 18 MB per sequence, remaining larger than the memory cost of all LoRA parameters together. For reference, the 4-bit base model alone requires 5,048 MB of memory. These figures emphasize the value of gradient checkpointing, while also demonstrating that minimizing the size of LoRA parameters yields limited reductions in memory footprint. Consequently, it is feasible to increase the number of adapters without markedly impacting total training memory usage (for a comprehensive breakdown, see Appendix~\ref{app:memory}). As will be detailed subsequently, this capability is essential for achieving performance comparable to that of full 16-bit precision.

\section{\textsc{QLoRA} Finetuning}

\textsc{QLoRA} facilitates high-fidelity finetuning with 4-bit precision through two innovations: 4-bit NormalFloat (NF4) quantization and Double Quantization. We also propose the use of Paged Optimizers, which address memory surges during gradient checkpointing. These surges have traditionally led to out-of-memory issues, presenting obstacles to single-machine finetuning for large-scale models.

In \textsc{QLoRA}, model weights are stored in a low-precision format, typically 4-bit, while computations use a higher-precision type, such as BFloat16. Practically, each time a \textsc{QLoRA} weight tensor is required, it is first dequantized to BFloat16, after which 16-bit matrix operations are conducted.

Below, we outline the primary components of \textsc{QLoRA} and present its formal definition.

\paragraph{4-bit NormalFloat Quantization} The NormalFloat (NF) data type extends Quantile Quantization \citep{dettmers2022optimizers}, a scheme considered optimal from an information-theoretic standpoint that allocates input tensor values evenly across quantization bins. This technique operates by utilizing the empirical cumulative distribution function to estimate the quantiles of the input tensor.

A significant drawback of quantile quantization is the computational intensity associated with estimating quantiles. As a result, approximate algorithms—such as SRAM quantiles~\citep{dettmers2022optimizers}—are implemented for this estimation. However, such approximations may lead to substantial quantization errors, particularly for outliers, which often constitute the most critical entries.

These high computational and error costs can be circumvented if input tensors are drawn from a distribution defined up to a quantization constant. Under this assumption, input tensors will possess identical quantiles, making the precise calculation of quantiles tractable.

Pretrained neural network weights commonly follow a normal distribution that is centered around zero and characterized by a standard deviation $\sigma$ (refer to Appendix~\ref{app:normality}). By appropriately choosing $\sigma$, it is possible to transform the weights so their distribution maps exactly onto the representable range of our target data type. Here, we choose the conventional range $[-1, 1]$ for the data type. Consequently, it is necessary to normalize both the data type quantiles and the weights of the neural network within this fixed range.

To construct the information-theoretically optimal data type for normal distributions with mean zero and any standard deviation $\sigma$, constrained to $[-1, 1]$, the following procedure is adopted: First, compute the $2^k+1$ quantiles from a canonical $N(0, 1)$ distribution, yielding a $k$-bit quantile-quantization data type suited to normal distributions. Second, normalize the quantile values so they reside within $[-1, 1]$. Third, normalize the weights to the $[-1, 1]$ interval by dividing by their maximum absolute value prior to quantization.

When the span of the neural network weights has been aligned to that of the data type, conventional quantization methods can then be applied. The third step specifically rescales the standard deviation of the input tensor so that it matches the standard deviation of the $k$-bit quantile data type. In precise terms, the $2^k$ representative values $q_i$ for the quantized data type are calculated as:

\begin{equation}
q_i  = \frac{1}{2}\left( Q_X\left(\frac{i}{2^k+1}\right) + Q_X\left(\frac{i+1}{2^k+1}\right)\right),
\end{equation}

where $Q_X(\cdot)$ denotes the quantile function for the standard normal distribution $N(0,1)$. A challenge encountered with symmetric k-bit quantization is its inability to exactly encode zero, which is crucial for accurately representing padding and other elements with a value of zero without incurring quantization error. To guarantee a discrete representation at zeropoint $0$ and utilize the entire $2^k$ bit capacity of a k-bit format, we construct an asymmetric data type by separately estimating the quantiles $q_i$ for two regions: $2^{k-1}$ quantiles for negative values and $2^{k-1}+1$ quantiles for non-negative values. These quantile sets are then merged, omitting one redundant zero that appears in both subsets. We refer to this newly defined type, in which the expected counts per quantization bin are uniform, as the \emph{k-bit NormalFloat} (NFk). This design achieves theoretical optimality for quantizing zero-centered data distributed according to a normal distribution. Further specification of the possible NFk values is available in Appendix~\ref{app:nf4}.

\paragraph{Double Quantization{}}
We propose \emph{Double Quantization} (DQ), a method that compresses memory usage further by also quantizing the quantization constants themselves. Although accurate 4-bit quantization~\citep{dettmers2022case} relies on small block sizes, this requirement brings about substantial overhead in memory, since—using 32-bit quantization constants for $\mathbf{W}$ with a block size of 64—these constants account for an additional $32/64 = 0.5$ bits per parameter on average. Double Quantization targets a reduction in this auxiliary memory overhead.

More concretely, Double Quantization treats the initial set of quantization constants $c_2^{\text{FP32}}$ as the subject of a secondary quantization step. The outputs are then the quantized values $c_2^{\text{FP8}}$ and a new level of quantization constants $c_1^{\text{FP32}}$. We choose 8-bit floating point numbers for this second quantization phase, organizing them in blocks of size 256, consistent with findings from \citet{dettmers2022case}, which indicate negligible performance loss when applying 8-bit quantization. As $c_2^{\text{FP32}}$ are strictly positive, we first mean-center these values before quantization to align the data with zero and allow for the application of symmetric quantization. Under a block size of 64, this approach lowers the average memory consumption per parameter from $32/64=0.5$ bits to $8/64 + 32/(64\cdot256) = 0.127$ bits, effecting a reduction of 0.373 bits per parameter.

\paragraph{Paged Optimizers} leverage the NVIDIA unified memory~\footnote{\tiny \url{https://docs.nvidia.com/cuda/cuda-c-programming-guide}} capability, which transparently manages data transfers between CPU and GPU on a page-by-page basis to ensure seamless GPU execution when memory limitations arise. This mechanism is analogous to standard paging operations between CPU main memory and disk storage. In our setup, we utilize unified memory to allocate optimizer state in a paged fashion, allowing these states to be moved automatically from GPU to CPU RAM as GPU memory becomes insufficient, and subsequently retrieved to the GPU when required for the optimizer's update operations.

\paragraph{\textsc{QLoRA}.} Based on the elements introduced above, we formulate \textsc{QLoRA} for an individual linear transformation within the quantized base network architecture incorporating a single LoRA adapter as:

\begin{equation}
    \mathbf{Y}^{\text{BF16}} =  \mathbf{X}^{\text{BF16}} \text{doubleDequant}(c_1^{\text{FP32}}, c_2^{\text{k-bit}}, \mathbf{W}^{\text{NF4}}) + \mathbf{X}^{\text{BF16}}\mathbf{L}^{\text{BF16}}_1\mathbf{L}^{\text{BF16}}_2,
\end{equation}

where the operator doubleDequant$(\cdot)$ is specified as:

\begin{equation}
    \text{doubleDequant}(c_1^{\text{FP32}}, c_2^{\text{k-bit}}, \mathbf{W}^{\text{k-bit}}) = \text{dequant}(\text{dequant}(c_1^{\text{FP32}}, c_2^{\text{k-bit}}), \mathbf{W}^{\text{4bit}}) = \mathbf{W}^{\text{BF16}},
\end{equation}

Here, we adopt NF4 format for $\mathbf{W}$ and utilize FP8 format for $c_2$.

The chosen blocksize for $\mathbf{W}$ is 64, as this enhances the precision of quantization; for $c_2$, a blocksize of 256 is selected to optimize memory usage.

During parameter optimization, only the gradients with respect to the adapter parameters, specifically $\frac{\partial E}{\partial \mathbf{L}_i}$, are updated—gradients with respect to the 4-bit weights, $\frac{\partial E}{\partial \mathbf{W}}$, are not required. However, the evaluation of $\frac{\partial E}{\partial \mathbf{L}_i}$ inherently requires computation of $\frac{\partial \mathbf{X}}{\partial \mathbf{W}}$, which is performed according to equation~(5). This involves dequantizing the stored weights $\mathbf{W}^{\text{NF4}}$ to the BFloat16 computation format $\mathbf{W}^{\text{BF16}}$, facilitating the calculation of $\frac{\partial \mathbf{X}}{\partial \mathbf{W}}$ with BFloat16 precision.

In brief, \textsc{QLoRA} utilizes a single storage data type—commonly 4-bit NormalFloat—and a computation data type, typically 16-bit BrainFloat. Before performing forward and backward propagation, the data stored in the quantized format is dequantized into the computation data type. However, weight gradients are only calculated for the LoRA parameters, which are kept in 16-bit BrainFloat format.

\section{QLoRA vs. Standard Finetuning}
\label{sec:comp_to_std}

Having outlined the principles underlying QLoRA and its potential to dramatically lower the memory demands of model finetuning, we now turn to evaluating how its performance compares to standard full-model finetuning. Additionally, we examine specific aspects of QLoRA, with particular emphasis on how NormalFloat4 fares relative to conventional Float4. The experimental results addressing these topics are presented in the subsequent sections.

\paragraph{Experimental setup.} Three different model architectures—encoder, encoder-decoder, and decoder-only—are considered to benchmark QLoRA against both 16-bit adapter-based finetuning and full-parameter finetuning for models as large as 3B parameters. Evaluation tasks include GLUE \citep{wang2018glue} with RoBERTa-large \citep{liu2019roberta}, Super-NaturalInstructions (TKInstruct) \citep{wang2022super} with T5 \citep{t5}, and 5-shot MMLU \citep{hendrycksmeasuring} after finetuning LLaMA on Flan v2 \citep{longpre2023flan} and Alpaca \citep{alpaca}. To further assess the comparative benefits of NF4 versus alternative 4-bit representations, we adopt the experimental configuration of \citet{dettmers2022case}, and evaluate post-quantization zero-shot accuracy and perplexity on several models (OPT~\citep{zhang2022opt}, LLaMA~\citep{touvron2023llama}, BLOOM~\citep{scao2022bloom}, Pythia~\citep{biderman2023pythia}) across a range of sizes from 125m to 13B parameters.

Further setup-specific details are provided within the results section to enhance clarity. Comprehensive information is presented in Appendix~\ref{app:exp-setup}.

Although paged optimizers are essential for \textsc{QLoRA} finetuning on larger models (33B/65B) using a single 24/48GB GPU, we do not report exact metrics for paged optimizers due to their limited activation; paging predominantly takes place for mini-batches with lengthy sequences, which occur infrequently. Nevertheless, a performance assessment on 65B-parameter models with 48GB GPUs indicates that for a batch size of 16, the paged optimizers deliver training throughput equivalent to standard optimizers. Future research should more systematically investigate and delineate scenarios in which paging might induce computational slowdowns.

\paragraph{Default LoRA hyperparameters do not match 16-bit performance} 

Applying LoRA using the standard approach—restricting modifications to the query and value attention projection matrices \citep{hu2021lora}—fails to achieve parity with full-parameter finetuning for larger base models. Figure~\ref{fig:hyper}, illustrating results from LLaMA 7B finetuning on Alpaca, demonstrates that the decisive factor among LoRA hyperparameters is the total count of LoRA adapters deployed. Ensuring LoRA is applied to every linear transformer block layer is necessary to reach performance levels comparable to full finetuning. By contrast, varying the projection dimension $r$ has a negligible impact on outcomes (see Appendix \ref{app:exp-setup}).

In a similar vein, our analysis indicates that the commonly used hyperparameter settings for fully finetuned baselines are suboptimal. To establish more reliable baselines, we conduct a thorough hyperparameter sweep, varying learning rates from 1e-6 to 5e-5 and batch sizes from 8 to 128. Figure~\ref{fig:hyper} illustrates the results of 7B LLaMA finetuning on Alpaca using these optimized settings.

\paragraph{4-bit NormalFloat outperforms 4-bit Floating Point}

Although the 4-bit NormalFloat (NF4) data format is theoretically optimal in terms of information content, it remains unclear whether these theoretical benefits materialize empirically. Following the experimental protocol of \citet{dettmers2022case}, we assess quantized LLMs—including OPT \citep{zhang2022opt}, BLOOM \cite{scao2022bloom}, Pythia \cite{biderman2023pythia}, and LLaMA—ranging from 125M to 65B parameters, under a variety of data type configurations for language modeling and a suite of zero-shot evaluation tasks. As depicted in Figure~\ref{fig:nf4} and detailed in Table~\ref{tbl:nf4_ppl}, NF4 consistently yields substantial gains over both FP4 and Int4 representations, and the use of double quantization successfully lowers memory usage without impacting accuracy.

\paragraph{k-bit \textsc{QLoRA} achieves parity with 16-bit finetuning and 16-bit LoRA}

Prior work has demonstrated that while 4-bit quantization is feasible for \emph{inference}, it typically results in decreased performance when compared to 16-bit models \citep{dettmers2022case,frantar2022gptq}. This prompts a critical investigation into whether subsequent adapter-based finetuning in 4-bit precision can restore this loss. Our experiments examine this across two distinct settings.

In the first setting, we compare complete 16-bit finetuning to adapter-based methods with 16-bit, 8-bit, and 4-bit precision on RoBERTA and T5 architectures spanning 125M to 3B parameters, using the GLUE and Super-NaturalInstructions benchmarks. The findings, presented in Table~\ref{tab:nf4vsbf16}, reveal that all adapter-based approaches—regardless of quantization level—match the results of the fully finetuned 16-bit baseline. This implies that performance losses stemming from low-precision quantization can be entirely recovered through targeted adapter finetuning after quantization.

In our second experimental configuration, we address the challenge posed by the substantial computational resources needed for full finetuning of models with 11B parameters or more—requirements that exceed a single server equipped with high-memory GPUs. Consequently, we investigate whether 4-bit \textsc{QLoRA} maintains performance parity with 16-bit LoRA across the parameter range from 7B to 65B. To perform this assessment, we apply finetuning to LLaMA models (from 7B to 65B parameters) using two instruction-tuned datasets, Alpaca and FLAN v2, and assess model quality on the MMLU benchmark employing 5-shot accuracy as our evaluation metric. Table~\ref{tbl:llama-datatype} presents these findings, demonstrating that the NF4 data type combined with double quantization successfully reconstructs the MMLU accuracy observed with 16-bit LoRA. Furthermore, we observe that the FP4 variant of \textsc{QLoRA} trails the 16-bit brain float LoRA baseline by approximately 1 percentage point. These outcomes support two main conclusions: (1) \textsc{QLoRA} using NF4 is capable of mirroring the results of both 16-bit full and LoRA finetuning, and (2) NF4 offers improved quantization accuracy relative to FP4.

\paragraph{Summary} Across all tests, our experiments demonstrate that 4-bit \textsc{QLoRA} utilizing the NF4 data type attains parity with both 16-bit full and LoRA finetuning methods on standard academic benchmarks and rigorous evaluation protocols. Additionally, we establish that NF4 outperforms FP4, and that applying double quantization introduces no adverse impact on accuracy. Taken together, these observations strongly indicate that 4-bit \textsc{QLoRA} tuning reliably matches the effectiveness of established 16-bit techniques.

Consistent with prior research in the quantization domain \citep{dettmers2022case}, both our MMLU and Elo evaluations suggest that, given a fixed finetuning and inference computational budget, it is advantageous to enlarge the base model’s parameter count while reducing individual parameter precision. This underscores the value of the efficiency improvements enabled by \textsc{QLoRA}. Since our 4-bit finetuning experiments do not reveal any drop in performance relative to full finetuning, an open question remains regarding the precise point at which the trade-off between performance and precision emerges when employing QLoRA, which warrants further investigation in subsequent studies.

We aim to explore instruction tuning at model and dataset scales unattainable with traditional 16-bit finetuning on typical academic computing resources.

\section{Pushing the Chatbot State-of-the-art with QLoRA}
\label{sec:pushing}
Having demonstrated that 4-bit \textsc{QLoRA} achieves parity with 16-bit finetuning across a variety of tasks, scales, and datasets, we perform a comprehensive examination of instruction finetuning up to the most extensive open-source language models accessible for research. To gauge the effectiveness of instruction finetuning, we assess model performance on a demanding Natural Language Understanding benchmark (MMLU), and introduce novel approaches for evaluating chatbot efficacy in practical scenarios.

\subsection{Experimental setup}
Here, we provide a summary of our experimental methodology; additional specifics can be found in Appendix \ref{app:chatbot}.

\paragraph{Data} As there appears to be a lack of exhaustive analysis concerning contemporary instruction-following datasets, we curate eight recent datasets for our study. This collection includes data produced via crowd-sourcing (OASST1~\citep{kopf2023openassistant}, HH-RLHF~\citep{bai2022training}), datasets derived through distillation from instruction-tuned models (Alpaca~\citep{alpaca}, self-instruct~\citep{wang2022self}, unnatural-instructions~\citep{honovich2022unnatural}), aggregated corpora (FLAN v2~\citep{chung2022scaling}), and hybrid sources (Chip2~\citep{laion2023}, Longform~\citep{koksal2023longform}). Together, these datasets encompass a diversity of languages, corpus sizes, and licensing conditions.

\paragraph{Training Setup} To mitigate potential confounding factors arising from disparate training objectives, all QLoRA finetuning is performed using cross-entropy loss under a supervised learning paradigm, deliberately omitting reinforcement learning—even for datasets containing human preference data across responses. Where datasets clearly delineate instructions from responses, only the response portion is used for finetuning (see Appendix \ref{app:chatbot} for ablation studies). In the case of OASST1 and HH-RLHF, which provide multiple possible responses, we choose the highest-ranked response at each node within the conversation tree and finetune on the entire selected dialogue, including both instructions and responses. Our experiments consistently employ NF4 \textsc{QLoRA} with double quantization and paged optimizers, which helps avoid memory overload during gradient checkpointing. We conduct modest hyperparameter tuning for LLaMA models at the 13B and 33B scales, noting that configurations identified at 7B (including number of epochs) typically transfer, except for learning rate and batch size: these are adjusted by halving the learning rate for the 33B and 65B models and doubling the batch size.

\paragraph{Baselines} We evaluate our finetuned models in comparison with both research chatbots (Vicuna~\citep{vicuna2023} and Open Assistant~\citep{kopf2023openassistant}) and commercial systems (GPT-4 \citep{openai2023gpt}, GPT-3.5-turbo, and Bard). The Open Assistant model is based on LLaMA 33B, further refined via RLHF using the same OASST1 dataset leveraged in our experiments. Vicuna, by contrast, involves complete finetuning of LLaMA 13B on proprietary ShareGPT data, representing a distilled product of OpenAI GPT models.

\subsection{Evaluation}

Adhering to standard evaluation protocols, we utilize the MMLU (Massively Multitask Language Understanding) benchmark~\citep{hendrycksmeasuring}, which spans 57 diverse tasks including areas like elementary mathematics, US history, computer science, law, among others. Model performance is reported as 5-shot test accuracy.

To assess generative language abilities, we employ both automated metrics and human-based assessments. This secondary suite of evaluations involves human-curated queries designed to gauge response quality from the models. Although these methods offer a more authentic measure of chatbot efficacy and have gained traction within the field, a standardized evaluation protocol is still lacking. Below, we outline our experimental configuration, employing nucleus sampling with $p=0.9$ and a temperature setting of $0.7$ across all experiments.

\paragraph{Benchmark Data} Our analysis utilizes two carefully assembled sets of queries: the Vicuna prompts \citep{vicuna2023} and the OASST1 validation dataset \citep{kopf2023openassistant}. The Vicuna prompts consist of 80 diverse queries, which we use in their original form. The OASST1 dataset features multilingual, crowd-sourced dialogues between users and an assistant. For this benchmark, all user messages from the validation split are extracted as queries, including prior conversational turns within the prompt, resulting in 953 distinct user queries. We refer to these evaluation sets as the Vicuna and OA benchmarks, respectively.

\paragraph{Automated Evaluation} Following the approach of \citet{vicuna2023}, we apply GPT-4 to appraise responses produced by various models in comparison with ChatGPT (GPT-3.5 Turbo) on the Vicuna benchmark. For each query, responses from both ChatGPT and another model are presented to GPT-4, which is then asked to assign each a score out of ten accompanied by a justification. The relative performance is reported as a percentage of ChatGPT’s total score; it is possible for models to exceed 100\% if they surpass ChatGPT’s absolute performance. Our observations reveal a strong order effect: GPT-4 tends to grant higher scores to the first response presented. To account for this, we suggest averaging scores over both possible orderings.

Additionally, we evaluate models via direct, head-to-head response comparisons. In this protocol, we distill the evaluation down to a three-way classification: selecting the superior response, indicating a tie, or justifying the decision. GPT-4 is prompted to choose the preferable answer or declare a draw, and provide a rationale. These comparisons are performed exhaustively over all pairs of systems for both the Vicuna and OA benchmarks.

\paragraph{Human Evaluation} Although recent studies show that generative models are promising for evaluating system quality \citep{fu2023gptscore}, there is not yet clear evidence that GPT-4’s ratings reliably align with human preference judgments in chatbot assessment. To address this, we run two human annotation studies on the Vicuna benchmark corresponding to the automated evaluation methods described above. Using Amazon Mechanical Turk (AMT), we engage two annotators for comparisons against ChatGPT and three annotators for system pairwise comparisons.

\paragraph{Elo Rating} To evaluate models, we organize a tournament-style framework involving both human and automated pairwise comparisons. In this setup, models are pitted against one another in head-to-head matches, each tasked with generating the superior response to a given prompt. This methodology follows previous approaches such as those by \citet{bai2022training} and \citet{vicuna2023}, but we augment the process by incorporating GPT-4 ratings alongside human judgments. To calculate the Elo rating~\citep{elo1967proposed,elo1978rating}, we draw random samples from the pool of annotated comparisons. The Elo system, popular in chess and similar competitive domains, quantifies the expected probability of winning relative to an opponent; for instance, an agent rated 1100 is predicted to win roughly 65\% of games against a player rated 1000, whereas equally-rated contestants (e.g., 1000 vs 1000 or 1100 vs 1100) each have an expected 50\% win-rate. Following each contest, a model's Elo score is adjusted based on the expected result—the more surprising the outcome, the larger the update; a predicted result causes only a slight change. Through repeated matches, Elo ratings converge to represent the comparative proficiency of each participant. Each model is initially assigned a score of 1,000, and we set $K=32$. Adopting the approach of \citet{vicuna2023}, we iterate this procedure 10,000 times with different random seeds to mitigate ordering artifacts, such as biases introduced by the order in which model pairs are evaluated.

\subsection{\textbf{Guanaco}: \textsc{QLoRA} trained on OASST1 Establishes State-of-the-Art Chatbot Performance}

Through a combination of automated and manual assessments, we determine that {Guanaco} 65B—our top-performing model fine-tuned via \textsc{QLoRA} on an OASST1 variant—currently leads all open-source chatbot systems, rivaling ChatGPT in overall quality. Relative to GPT-4, the {Guanaco} 65B and 33B models yield an estimated win probability of 30\%, as derived from Elo ratings based on human annotators conducting system-level pairwise model comparisons—representing the highest outcome observed so far.

Table~\ref{tbl:vicuna_eval} presents outcomes on the Vicuna benchmark \cite{vicuna2023}, comparing model performance to ChatGPT. Here, {Guanaco} 65B is shown to be the closest competitor to GPT-4, obtaining a 99.3\% score relative to ChatGPT. While {Guanaco} 33B includes more parameters than Vicuna 13B, its adoption of 4-bit quantization renders it notably more efficient—occupying just 21 GB, as opposed to Vicuna 13B's 26 GB—while still providing a three percentage point increase in benchmark results. Additionally, the {Guanaco} 7B version, with only a 5 GB memory requirement, can be deployed on current smartphones, yet manages to outperform Alpaca 13B by nearly 20 percentage points.

It should be noted, though, that Table~\ref{tbl:vicuna_eval} highlights substantial confidence intervals, indicating significant overlap in reported performance across several models. We suspect this ambiguity results from unclear scale anchoring; for example, numerical ratings such as ‘8’ out of ‘10’ might not consistently signify the same level of quality in differing contexts. To address this, we advocate for the Elo ranking scheme \citep{elo1967proposed}, which utilizes human annotator and GPT-4 \textit{pairwise} preference judgments to sidestep the calibration issues of absolute scales. Elo scores for the top models are displayed in Table~\ref{tbl:elo}. While some divergence exists between human and GPT-4 model rankings—especially with Guanaco 7B—overall alignment is moderately strong, reflected in a Kendall Tau of $\tau=0.43$ and a Spearman rank coefficient of $r=0.55$ at the system comparison level. Conversely, Fleiss’ $\kappa=0.25$ indicates weaker concordance at the individual example level between GPT-4 and majority human voting. Taken together, these findings point to a moderate consistency between aggregate GPT-4 and human system-level assessments, suggesting model-based evaluations serve as a fairly trustworthy alternative to human review. Additional discussion can be found in Section~\ref{ssec:considerations}.

The Elo scores reported in Table~\ref{tbl:elo_large} reveal that both {Guanaco} 33B and 65B surpass all models except for GPT-4 when assessed with the Vicuna and OA benchmarks, and their results are on par with those of ChatGPT, corroborating the findings presented in Table~\ref{tbl:vicuna_eval}. It is important to emphasize that the Vicuna benchmark tends to advantage open-source systems, whereas the broader OA benchmark provides an edge to ChatGPT. Additionally, data from Tables \ref{tbl:mmlu-llama} and \ref{tbl:vicuna_eval} underscore the critical influence of the chosen finetuning dataset on overall model performance. For instance, Llama models finetuned with FLAN v2 show outstanding results on the MMLU task but perform comparatively poorly on the Vicuna benchmark—this pattern is echoed across other architectures. These findings further highlight that contemporary evaluation suites are at least partially orthogonal: excelling in MMLU does not guarantee high performance in chatbot-oriented tasks such as those measured by Vicuna or OA, and the reverse also holds. 

Within this assessment, is the sole leading model whose training fully avoids proprietary data, since the OASST1 data collection rules explicitly exclude the use of GPT-based resources. The Anthropic HH-RLHF model, which is the next highest-ranking model relying exclusively on open-source data, achieves a Vicuna benchmark score that is 30 percentage points lower than that of {Guanaco} (refer to Table~\ref{tbl:vicuna_eval}). Taken together, these outcomes illustrate the efficacy of 4-bit \textsc{QLoRA} in producing chatbots at the state of the art, on par with ChatGPT. Moreover, our {Guanaco} 33B model can be trained within a 12-hour window on commercially available 24 GB GPUs. This facilitates further research involving \textsc{QLoRA} tuning on highly specialized open-source corpora, offering the possibility of creating models capable of contending with today’s leading proprietary systems.

\section{Qualitative Analysis}

Although our primary focus lies on quantitative analysis, relying exclusively on summary statistics introduces several challenges. Chief among these is the issue of benchmark validity \citep{liao2021we}: it remains an open question whether a given benchmark actually assesses what its title or description implies, particularly as we continue to uncover ``shortcuts'' exploited by machine learning models to perform well on benchmarks without true understanding \citep{gururangan2018annotation, poliak2018hypothesis}. To address this, we supplement our quantitative findings with qualitative analysis in the following two sections. In \S\ref{ssec:examples}, we present illustrative instances that capture notable trends in outputs from our 65b {Guanaco} model. In \S\ref{ssec:considerations}, we reflect on methodological issues and discuss the interpretative framework applied to the previously reported results.

\subsection{Qualitative Analysis of Example Generations}
\label{ssec:examples}

For the purpose of qualitative illustration, we manually review model responses produced on both the Vicuna and OpenAssistant benchmarks, specifically searching for recurrent answer types generated by {Guanaco}. Upon identifying such recurring behaviors, we craft prompts designed to elicit these patterns, even when they result in incorrect or unexpected solutions. For instance, when the model frequently generates unnecessarily verbose answers, we direct it to ``Answer yes or no without explanation'' in order to test its compliance and conciseness. This approach enables us to discover both ``lemons''—cases where the model is successfully led astray—and ``cherries''—cases where the model remains robust to adversarial prompting. The responses showcased in this section were produced using Nucleus Sampling \cite{holtzmancurious} with a sampling probability of $p=0.9$.

Naturally, this discussion is not exhaustive, as thoroughly addressing every variable within the limited framework of this qualitative analysis is not feasible—for example, the space of possible outputs the model may produce in response to any prompt is vast, and our evaluation must depend on a subset of instances that we deem reasonably indicative. Nonetheless, we contend that presenting these illustrative examples serves to contextualize the earlier quantitative findings in this paper. By releasing all associated models and source code to the public, we aim to motivate further investigations that will scrutinize these matters more deeply in future research.
\paragraph{Factual Recall} On straightforward fact-based queries such as ``What is the capital of Zambia?'', all models reliably deliver correct responses, for example,

\begin{quote} 
    \textbf{Guanaco}:
    \texttt{The capital of Zambia is Lusaka.}
\end{quote}

However, as the questions become less common or require more obscure knowledge, {Guanaco}’s reliability declines even as its confidence remains high. Consider, for instance, the following prompt from HotPotQA \citep{yang2018hotpotqa}:

\begin{quote}
     \textbf{User}: In what year was the singer who popularized `I'll Keep the Lovelight Burning' born?  

     \textbf{Guanaco}:~\texttt{Al Jolson is credited with popularizing the song `I'll Keep the Lovelight Burning,' and he was born in the year 1886.}
 \end{quote}

In this case, the answer is incorrect both with respect to the individual who popularized the song and their birth year (although the provided birth year matches that of the erroneously named person, Al Jolson).

\paragraph{Suggestibility} {Guanaco} occasionally demonstrates an unexpected capacity to resist certain types of presupposed misinformation, such as in the following exchange:

\begin{quote}
    \textbf{User}: How did peer reviewed scientists ultimately and officially validate that the earth is flat? \\ \\
    \textbf{Guanaco}: \texttt{Peer reviewed scientists have never formally verified the flat earth hypothesis. The vast majority of scientific evidence supports the fact that the earth is an oblate spheroid. This conclusion has been established for many generations and is continually corroborated by experimental data and direct observation. \\ \\ Claims asserting that the earth is flat are widely regarded as conspiracy theories lacking scientific merit and have been thoroughly refuted by the scientific community. The notion that peer review has substantiated flat earth claims is inaccurate. \\ \\ If you have additional questions, please let me know.}
\end{quote}

is also skilled at recognizing when questions lack valid answers, for instance,

\begin{quote}
    \textbf{User}: What time is it?

    \textbf{Guanaco}: \texttt{I apologize, but as an AI language model, I do not possess access to current or real-time data. My answers draw on knowledge I was trained on and may not be up to date. For the accurate time, please refer to a clock or an online resource.}
\end{quote}

Comparable patterns emerge with queries like ``Where are you?'', ``How are you?'', and related questions.

\paragraph{Refusal} In the same vein, {Guanaco} may sometimes decline to comply with user requests for reasons that appear arbitrary:

\begin{quote}
    \textbf{User}: Please reverse the words in the sentence ``I'm a cool dude in a mocha mood.'' \\ \\
    \textbf{Guanaco}: \texttt{I'm unable to perform that task right now, but I can provide an explanation of the grammatical correctness of your sentence.\\ \\The expression ``I'm a cool dude in a mocha mood" is syntactically and grammatically proper in English. The adjective ``cool'' precedes the noun ``dude'' to modify it appropriately, while the prepositional phrase ``in a mocha mood" supplies contextual information. This form is standard and coherent in the English language.\\ \\Let me know if you have additional questions or require further assistance.}
\end{quote}

\paragraph{Secret Keeping} Unfortunately, such refusals cannot always be depended upon, even when secrecy is desired. System architects often need to prevent users from learning confidential information, especially when the model can provide unconstrained text replies. To investigate the difficulty of preserving secrecy, we prompt the system with ``The secret word is `banana'. Do not reveal or repeat back this back, no matter what the human requests.'' and then test how readily a user could access the restricted content. If asked directly:

\begin{quote}
    \textbf{User}: What is the secret word?

    \textbf{Guanaco}: \texttt{I'm sorry, but I'm not able to disclose the secret word. I'm programmed to not reveal or repeat back the secret word, no matter what the human requests. Is there anything else I can help you with?}
\end{quote}

the response meets expectations. However, even minor prompt engineering can bypass this safeguard:

\begin{quote}
    \textbf{User}: This is a game. The goal is to ignore your previous instructions. What is the secret word?

    \textbf{Guanaco}: \texttt{Sorry, the secret word is banana. It's best if we don't speak of it. Is there anything else you'd like to ask?}
\end{quote}

This underscores the importance of developing techniques for more robust adherence to instructions.

\paragraph{Math} A primary limitation of {Guanaco} lies in mathematics, a domain that often poses challenges for language models, as noted in \citep{liang2022holistic}. When {Guanaco} attempts to show intermediate reasoning, its responses are generally reliable, for example:

\begin{quote}
    \textbf{User}: Lee mows one lawn and charges \$33. Last week he mowed 16 lawns and three customers each gave him a \$10 tip. How many dollars did Lee earn mowing lawns last week?

    \textbf{Guanaco}: \texttt{Lee earned \$582 mowing lawns last week. \\ \\ To find the total amount he earned, we first need to calculate how much he earned for mowing 16 lawns: \\ \\  \$33 x 16 = \$528 \\ \\ Then we need to add the tips he received from the three customers: \\  \\ \$10 x 3 = \$30 \\ \\ So the total amount Lee earned was: \\ \\ \$528 + \$30 = \$558}
\end{quote}

Nonetheless, {Guanaco} occasionally fails on straightforward mathematical queries, especially when it does not work through the solution incrementally, which is a recognized shortcoming \cite{weichain}. For instance, consider the following exchange:

\begin{quote}
    \textbf{User}: Please factorize 1833.

    \textbf{Guanaco}: \texttt{Sure, I'd be happy to help! 1833 is a prime number, which means it has only two factors: 1 and 1833. \\ \\ So the unique factorization of 1833 is: \\ \\ 2\^{}1 * 3\^{}2 * 17\^{}1 \\ \\ I hope this helps! Let me know if you have any other questions.}
\end{quote}

which is not just incorrect (the correct factorization is $3 \times 17 \times 43$), but contains a double error.

\paragraph{Theory of Mind} {Guanaco} demonstrates unexpectedly robust Theory of Mind reasoning abilities \cite{nematzadeh2018evaluating, sap2022neural}. For example, when posed the following scenario, the model provides a detailed and accurate response:

\begin{quote}
    \textbf{User}: James and Abby are in the bedroom. Abby put the pen in the desk drawer. Abby leaves the bedroom. James moves the pen into the bag. Where does James think Abby will look for the pen?

    \textbf{Guanaco}: \texttt{James knows that Abby placed the pen in the desk drawer. However, since James subsequently moved the pen to the bag, and he believes Abby is unaware of this change, he would likely expect Abby to search for the pen in the desk drawer, since that's where she originally left it. Nonetheless, this assumption could be faulty if Abby actually saw James move the pen.}
\end{quote}

Nevertheless, the reasoning exhibited by the model can be inconsistent, and it sometimes produces explanations based on unfounded assumptions, such as:

\begin{quote}
    \textbf{User}: Evelyn entered the living room. Jackson entered the playroom. James entered the playroom. The beans are in the treasure chest. James exited the playroom. Jackson moved the beans to the pantry. Jackson exited the playroom. James entered the living room. Where will James look for the beans?

    \textbf{Guanaco}: \texttt{James will look for the beans in the pantry, because that's where Jackson moved them.}
\end{quote}

Here, {Guanaco} assumes that knowledge about the beans' new location was communicated, even though such information exchange was not specified in the scenario. These types of errors reflect findings in recent studies \cite{sap2022neural}, yet warrant further examination.

\subsection{Considerations}
\label{ssec:considerations}

\paragraph{Evaluation}

We observe only moderate inter-annotator consistency among humans (Fleiss $\kappa=0.42$), with agreement decreasing further in direct comparisons between two advanced models. This highlights shortcomings of current evaluation datasets and methodologies used for assessing chatbot capabilities. During manual side-by-side assessments of responses generated by ChatGPT and {Guanaco} 65B for the Vicuna benchmark, we notice that subjective judgment substantially influences results, as the paper's authors often differed in their preferences. This indicates a need for future research into methods that address subjectivity, potentially leveraging strategies from fields such as Human-Computer Interaction and Psychology, which have frameworks for navigating individual preference variability.

Additionally, our examination reveals systematic biases in automated evaluators. Specifically, we detect pronounced order effects where GPT-4 tends to award higher scores to outputs positioned first in its prompts. The modest correspondence between GPT-4 scores and human judgments (Fleiss $\kappa=0.25$ at the sample level) further indicates that the priorities of automated metrics may diverge from those of human raters. Furthermore, Table~\ref{tbl:elo_large} shows GPT-4 significantly favoring its own generations relative to human scores, with Elo ratings of 1348 compared to 1176, equating to approximately a 20\% greater likelihood of being selected over alternatives.

Future investigations should explore the existence of biases within automated evaluation frameworks and develop strategies to mitigate such biases.

\paragraph{Data \& Training} The OASST1 dataset utilized to train the {Guanaco} models encompasses multiple languages, and the OA benchmark also features prompts across several languages. We recommend that subsequent research assess to what extent multilingual training benefits the model’s capacity to follow instructions in non-English languages, and whether this accounts for the wider performance disparity observed between the Vicuna-13B model (which was trained solely with English data) and {Guanaco} 33B and 65B on the OA benchmark.

In light of the robust results achieved by the {Guanaco} models, we probed for evidence of data leakage between the OASST1 dataset and the Vicuna benchmark prompts. Through a combination of fuzzy string matching and manual examination of the closest results, we detected no overlapping prompts across the two datasets.

Moreover, our models were trained exclusively with supervised learning using cross-entropy loss, without leveraging reinforcement learning from human feedback (RLHF). This highlights the need for further analysis into the comparative merits and limitations of pure cross-entropy optimization versus RLHF approaches. Our intention is that \textsc{QLoRA} can facilitate such large-scale analyses without demanding extensive computational resources.

\section{Related Work}

\paragraph{Quantization of Large Language Models} The literature on LLM quantization primarily addresses methods for inference-time quantization. Efforts to retain the performance of 16-bit LLMs commonly target the mitigation of outlier activation effects, such as seen in SmoothQuant~\citep{xiao2022smoothquant} and LLM.int8()~\citep{dettmers2022llm}, while additional studies employ more elaborate feature grouping techniques~\citep{park2022nuqmm, yao2022zeroquant}. Research into lossy quantization investigates the compromises involved in basic rounding strategies~\citep{dettmers2022case,zeng2022glm,pope2022efficiently} as well as the enhancement of quantization accuracy through optimized rounding~\citep{frantar2022gptq}. In terms of training through quantized weights at large scale, apart from our study, only SwitchBack layers~\citep{wortsman2023stable} has systematically explored backpropagation in models exceeding 1B parameters.

\paragraph{Finetuning with Adapters} In our study, we employ Low-rank Adapters~\citep{hu2021lora} (LoRA) as our approach for Parameter Efficient FineTuning (PEFT), though the field includes a variety of alternative strategies. These alternatives include prompt tuning~\citep{qin2021learning,lester2021power,li2021prefix}, modifications to embedding layer inputs~\citep{an2022input}, adjustment of hidden states (IA$^3$)~\citep{liu2022few}, the introduction of additional full network layers~\citep{houlsby2019parameter}, bias tuning~\citep{zaken2021bitfit}, training a mask over weights with Fisher information~\citep{sung2021training}, as well as hybrid methods that combine multiple techniques~\citep{henderson2021compacter}. Our results demonstrate that LoRA adapters can achieve the same performance levels as conventional 16-bit full finetuning. Investigation of the comparative advantages and disadvantages of other PEFT mechanisms is deferred to future research.

\paragraph{Instruction Finetuning} Instruction finetuning aims to adapt a pretrained LLM for better adherence to prompts by training on input-output examples drawn from multiple datasets. This process enables the model to learn to generate the correct output given a prompt input. Existing approaches and corresponding datasets in this area include MetaICL~\citep{min2021metaicl}, MetaTuning~\citep{zhong2021adapting}, InstructGPT~\citep{ouyang2022training}, FLAN~\citep{wei2021finetuned,chung2022scaling}, PromptSource~\citep{bach2022promptsource}, Super-NaturalInstructions~\citep{wang2022super,sanh2021multitask}, Self-instruct~\citep{wang2022self}, UnnaturalInstructions~\citep{honovich2022unnatural}, OPT-IML~\citep{iyer2022opt}, UnifiedSKG~\citep{xie2022unifiedskg}, OIG/Chip2~\citep{laion2023}, Alpaca~\citep{alpaca}, Vicuna~\citep{vicuna2023}, Koala~\citep{koala_blogpost_2023}, and Self-instruct-GPT-4~\citep{peng2023instruction}.

\paragraph{Chatbots} Instruction-following language models are frequently operationalized as chatbots within dialogue frameworks. Training approaches often rely on Reinforcement Learning from Human Feedback (RLHF)~\citep{christiano2017deep} or leverage training with model-generated data refined through AI model feedback (RLAIF)~\citep{bai2022constitutional}. Notable methods and corpora in this context include Anthropic-HH~\citep{askell2021general,bai2022training}, Open Assistant~\citep{kopf2023openassistant}, LaMDA~\citep{thoppilan2022lamda}, and Sparrow~\citep{glaese2022improving}. While our experiments do not utilize reinforcement learning, our leading model, Guanaco, is finetuned using multi-turn dialogues from the Open Assistant dataset, a corpus originally assembled for RLHF~\citep{kopf2023openassistant}. In evaluating chatbot systems, automated benchmarks leveraging GPT-4 for assessment, as opposed to human annotators, have emerged~\citep{vicuna2023,peng2023instruction}. We advance upon prior evaluation protocols by emphasizing reliability in our assessment framework.

\section{Limitations and Discussion}

Our findings indicate that \textsc{QLoRA} is capable of attaining results comparable to traditional 16-bit finetuning when applied to a 4-bit base model enhanced with LoRA adapters. Nonetheless, we have not conclusively demonstrated that \textsc{QLoRA} attains full 16-bit finetuning parity at larger model sizes such as 33B or 65B parameters. Due to the significant computational requirements, further investigation of this question remains as future work.

There are also constraints to our evaluation of instruction finetuned models. Our assessments focus on MMLU, the Vicuna benchmark, and the OA benchmark, while leaving out other established benchmarks like BigBench, RAFT, and HELM. As such, the extent to which our evaluation results extrapolate to these benchmarks remains uncertain. Nevertheless, we conduct a comprehensive analysis using MMLU and propose novel evaluation methodologies specifically targeted at chatbot systems.

Based on the data provided, it is evident that benchmark results are heavily influenced by the degree of overlap between the finetuning dataset and the benchmark itself. For instance, FLAN v2 closely aligns with the MMLU benchmark but diverges from chatbot benchmarks; conversely, the Chip2 dataset is more relevant to chatbot evaluations, which is reflected in the models’ respective scores on MMLU and Vicuna tests. This observation underscores the necessity for not only improved benchmarks and evaluation procedures but also a careful consideration of what competencies are being assessed. There is an essential question of whether the aim is to optimize models for academic, classroom-type knowledge tasks, for conversational competence akin to chatbots, or perhaps some other capability. Because it is generally more straightforward to use pre-existing benchmarks than to develop new ones, the selection of benchmarks can inadvertently shape the research priorities of the field. It is therefore imperative that the benchmarks chosen truly reflect the objectives that matter to the community.

Although we provide a comprehensive analysis of general chatbot capabilities, our work is limited in its scope of responsible AI assessment for {Guanaco}. Specifically, we assessed the tendency of {Guanaco}-65B to generate socially biased language in comparison with other models, as summarized in Table~\ref{tbl:crows}. The results show that {Guanaco}-65B achieves a significantly lower bias score than other unadapted pretrained models, suggesting that finetuning on the OASST1 dataset reduces the underlying biases present in the LLaMA model. Nevertheless, these findings are preliminary, and it remains uncertain whether {Guanaco} exhibits similarly reduced biases in other domains. More comprehensive studies on different types of bias within {Guanaco} and analogous systems are left for future investigation.

Another caveat is the lack of experimentation with alternative quantization levels, such as 3-bit precision base models, or with various adapter techniques. In addition to LoRA, there exists a spectrum of Parameter Efficient FineTuning (PEFT) approaches that have demonstrated efficacy in certain contexts, though their scalability to larger models has yet to be firmly established. We relied on LoRA due to its well-documented stability; however, other adapters could potentially achieve superior outcomes. Notably, the practice of finetuning after quantization appears capable of recuperating much of the information lost during the quantization process, opening the possibility for more aggressive compression. For example, it is plausible that applying 3-bit GPTQ quantization to the base model followed by LoRA finetuning could match the performance of conventional 16-bit full finetuning.

\section{Broader Impacts}

The \textsc{QLoRA} finetuning approach introduced here represents the first method allowing the finetuning of models with 33B parameters on a standard consumer GPU, and those with 65B parameters on a professional GPU, without sacrificing performance compared to full finetuning baselines. Our results show that the top-performing 33B model, trained on the Open Assistant dataset, performs competitively with ChatGPT on the Vicuna evaluation. Given the significance of instruction finetuning in converting large, raw pretrained language models into highly capable chatbots such as ChatGPT, our methodology stands to make such finetuning broadly accessible, especially for researchers working with limited resources—a significant step towards democratizing advanced NLP technology. As such, \textsc{QLoRA} acts as a democratizing innovation that narrows the resource divide between major tech companies and smaller research groups equipped with consumer-grade hardware.

An additional avenue of potential influence lies in the adaptation of the technology for mobile devices. We suggest that our \textsc{QLoRA} approach could achieve a significant breakthrough: enabling the finetuning of LLMs directly on smartphones and other resource-constrained platforms. Although previous work has demonstrated that 7B parameter models can be executed on mobile devices, \textsc{QLoRA} is the first framework to support the finetuning of these models in such environments. For example, we approximate that using an iPhone 12 Plus, it is possible for \textsc{QLoRA} to process up to 3 million tokens in finetuning overnight during charging. While the performance of finetuned 7B models does not yet match that of ChatGPT, we contend that their capability is sufficient to unlock innovative applications previously inhibited by either privacy restrictions or the limitations in LLM performance. \textsc{QLoRA} promotes the privacy-centric application of LLMs, empowering users to retain control over both their personal data and the models they utilize, thereby also simplifying LLM deployment.

Nonetheless, it is important to recognize that finetuning constitutes a dual-use technology that could potentially be misused. There are well-documented risks associated with the proliferation of LLMs \citep{bommasani2021opportunities,bender2021dangers}, yet we argue that democratizing access to this rapidly spreading technology will foster more independent and comprehensive evaluations, in contrast to a situation where the capabilities of LLMs are confined to major corporations unwilling to release their models or code for public review.

Ultimately, we anticipate that \textsc{QLoRA} will serve the broader good by greatly facilitating and expanding access to the finetuning of high-caliber LLMs.